\section{Introduction}

La convergence de la géométrie algébrique et de l'information quantique a ouvert de nouvelles voies pour comprendre les phénomènes macroscopiques. Plus particulièrement, la relation entre les fausses fonctions modulaires de Ramanujan, les décomptes d'états BPS des compactifications de surfaces K3, et l'émergence de la matière noire floue macroscopique (FDM) constitue le cœur de notre cadre à Double Échelle (Dual-Scale). Dans cet article, nous consolidons les fondements mathématiques de cette relation, vérifiés de manière computationnelle via Lean 4, et nous démontrons sa cartographie native dans les réseaux tensoriels holographiques.

En contournant les approches classiques de la géométrie différentielle et en utilisant une arithmétique rationnelle exacte, nous établissons un pont rigide depuis les états quantiques microscopiques des cordes (UV) jusqu'à la physique macroscopique observable (IR).
